% !TeX root = ../main.tex
\section{Introduction}
\label{sec:introduction}

Future flapping-wing aerial vehicles are expected to do more than maintain stable flight. A practical vehicle may need to regulate flight while also using perception, interacting with targets, grasping or perching, and making task-level decisions in outdoor environments. This broader autonomy goal makes reinforcement learning attractive, because flight states, visual observations, contact or interaction cues, and task rewards can be handled within a single training framework~\cite{kober2013rlsurvey}. The same breadth also makes direct hardware training difficult. Real flight tests expose the vehicle to actuator limits, wind, impact, and recovery failures, and broad generalization would require data from many successful and unsuccessful operating conditions.

Real flight logs are therefore valuable but not sufficient by themselves. They contain the actual vehicle response, actuator behavior, sensing effects, and outdoor disturbances, yet they usually cover only the state distribution that has already been flown. For a bird-scale flapping-wing platform, collecting enough logs to cover aggressive maneuvers, rare disturbances, failure cases, and wide operating envelopes can be impractical. This limitation pushes future learning, controller design, trajectory replay, and task-level evaluation toward simulation. A useful simulator, however, must contain a model of the flapping-wing dynamics that is both accurate enough to be meaningful and fast enough for repeated rollouts.

For flapping-wing vehicles, the main modeling bottleneck is aerodynamics. The vehicle generates forces and moments through periodic wing motion, wing flexibility, tail and body aerodynamics, and coupled motion between the fuselage and wings. High-fidelity computational fluid dynamics can resolve more of this flow physics, but it is too slow for large-scale controller training or repeated evaluation. Vortex-based and reduced-order approaches provide intermediate fidelity. Empirical and semi-empirical formulations sacrifice some physical detail for speed, but their structure makes them practical simulator priors. In this paper, DeLaurier's flapping-wing aerodynamic formulation~\cite{delaurier1993} is used as such a prior.

Speed alone does not make an empirical prior flight-realistic. A nominal aerodynamic calculation cannot fully capture airframe-specific geometry, wing deformation, actuator timing, tail effectiveness, body--wing interaction, sensor alignment, and state-estimation errors. Outdoor flight also adds wind and airframe-condition variation. Domain randomization and dynamics randomization can improve robustness when policies are transferred from simulation to hardware~\cite{tobin2017domainrandomization,peng2018dynamicsrandomization}, but randomization is most useful when the simulator prior and its systematic errors are understood. If a repeatable part of the gap is aerodynamic, the first engineering problem is to diagnose where the simulator prior differs from real flight and whether that residual can be predicted from onboard measurements.

Flight data have long been used to make flapping-wing models useful beyond idealized aerodynamic analysis. Early DelFly studies reconstructed aerodynamic forces and moments from motion-capture flight tests and identified simple linear or time-varying models, often at a flap-averaged or low-order modeling level~\cite{caetano_linear_2013,armanini_time-varying_2016}. Related grey-box and onboard-sensor-based methods have been reported for tailless flapping-wing micro air vehicles and longitudinal stability analysis~\cite{nijboer_longitudinal_2020,aurecianus_longitudinal_2022}. At larger scales, flight-data studies of ornithopters have examined force estimation, trim behavior, and periodic state variation in forward flight~\cite{aminietal2020largescale}. More recent avian-inspired platforms show that real flight data can expose flapping-period structure: phase-functioned neural models have predicted equivalent aerodynamic forces and moments from flight data~\cite{zhao2025phase}, and learned periodic corrections from onboard inertial and GNSS data have improved state estimation for a $1.7~\mathrm{m}$ wingspan vehicle~\cite{qian_learning_2026}.

Controlled aerodynamic datasets provide a complementary route to data-driven flapping-wing models. Neural models have been used to predict unsteady flapping-airfoil loads~\cite{kurtulus_ability_2009}, transient aerodynamic properties of flapping-wing systems~\cite{lan_neural_2022}, and robust reduced models for hovering flapping wings~\cite{calado_robust_2023}. Recent surrogate models also predict time-history lift, thrust, and moment curves for optimization using computational or experimental datasets~\cite{wang_time-history_2024,corban_discovering_2023,chen_data-driven_2026}. These studies demonstrate the value of data-driven aerodynamic modeling, but many rely on prescribed kinematics, wind-tunnel measurements, computational data, or bench-top experiments. Learning methods have also improved flapping-wing control, estimation, and simulation~\cite{he_adaptive_2017,qian_neural_2023,qian_learning_2023,fei_flappy_2019}. The remaining need for simulation is to connect a fast, structured aerodynamic prior to instantaneous six-axis effective force and moment labels reconstructed from outdoor flight logs of the complete vehicle.

Our platform provides the measurements needed for that connection. The vehicle flies with a PX4-based fixed-wing stack~\cite{meier2015px4} while recording synchronized state estimates, airdata, actuator commands, flapping phase, and flapping frequency. These logs allow the body-frame effective wrench required to explain the measured rigid-body motion to be reconstructed and compared directly with the wrench predicted by a DeLaurier-style simulator model. The comparison turns real flight data into a diagnostic for the aerodynamic component of the simulator gap.

This paper presents a real-flight effective-wrench residual modeling pipeline for a bird-scale flapping-wing vehicle. The paper does not train a reinforcement-learning controller. Instead, it develops one engineering step toward real-flight-corrected simulation for future learning and control. In this work, body-frame six-axis effective wrench labels are reconstructed from outdoor flight logs, aligned with a DeLaurier-style aerodynamic simulator prior, and used to calibrate that prior using training logs only. A causal neural residual model then predicts the remaining discrepancy from onboard-measurable variables. Finally, the residual is analyzed by wingbeat phase, flapping frequency, and flight condition to determine whether the simulator mismatch is structured rather than random. The main contributions are:
\begin{enumerate}
  \item We construct real-flight six-axis effective force and moment labels from outdoor flight logs of a bird-scale flapping-wing vehicle and align them with a DeLaurier-style aerodynamic simulator prior to define an effective-wrench residual modeling problem.
  \item We combine a calibrated aerodynamic prior with a causal neural residual estimator that predicts the remaining simulator-to-flight discrepancy from onboard measurements. The resulting hybrid model is evaluated on held-out logs and date-stratified whole-log cross-validation.
  \item We use phase-binned, condition-binned, and frequency-domain diagnostics to interpret the learned residual as an engineering diagnostic for simulator mismatch. The results show that the calibrated aerodynamic prior leaves repeatable wingbeat-synchronous and flight-condition-dependent force residuals, and that the neural residual estimator removes much of this structured discrepancy.
\end{enumerate}

The remainder of the paper is organized as follows. Section~\ref{sec:dataset} describes the flight platform, logging setup, and effective-wrench reconstruction. Section~\ref{sec:method} presents the calibrated aerodynamic prior and causal residual model. Section~\ref{sec:evaluation} defines the evaluation protocol, and Section~\ref{sec:results} reports prediction performance and residual-structure analysis.
